% GENERATED FILE. Edit canonical lessons, then run npm run generate:books.
% canonical-source-hash: fnv1a64:5ea0a5f3eb773784
% canonical-lessons: MR-C415-suvidha, MR-C415-bichana, MR-C415-ghongdi, MR-C415-antharne, MR-C415-pangharne

\chapter{Making the Bed}
\label{ch:mr-making-the-bed}
\hlchaptermodality{415}

\begin{chapteropening}
\textbf{By the end of this chapter:} \emph{I can talk about amenities, bedding and blankets, and making up a bed.}

Complete the last lesson of chapter 415: I can talk about amenities, bedding and blankets, and making up a bed.
\end{chapteropening}

\section[\texorpdfstring{\mr{सुविधा} (suvidhā)}{सुविधा (suvidhā)}]{\mr{सुविधा} (suvidhā) — an amenity}

\label{lesson:MR-C415-suvidha}



\begin{quote}
Before the new one: say the Marathi for to ask after, then the Marathi for to drop by.
\end{quote}

\subsection*{You'll want to know: \mr{सुविधा}}
\textbf{\mr{सुविधा}} — \emph{suvidhā} — \textquotedblleft{}an amenity\textquotedblright{}.

\begin{quote}
\mr{या} \mr{विश्रामगृहात} \mr{सगळ्या} \mr{सुविधा} \mr{आहेत}.
\end{quote}

\begin{quote}
\emph{yā viśrāmgṛhāt sagaḷyā suvidhā āhet.}
\end{quote}

\begin{quote}
This guest house has all the amenities.
\end{quote}

Feminine.

A feminine noun; the plural is the same in the direct case.

\begin{cousinweb}
\textbf{\mr{सुविधा}} (\emph{suvidhā}) is a modern word built from Sanskrit \emph{su-}, \textquotedblleft{}good\textquotedblright{}, and \emph{vidhā}, \textquotedblleft{}arrangement, manner\textquotedblright{} (Monier-Williams).
\end{cousinweb}

\begin{grammarlens}[title={What you've built}]
One more word for getting around.
\end{grammarlens}

\subsection*{Your turn}
\begin{itemize}
\raggedright
  \item \emph{Say it:} \emph{suvidhā}
  \item \emph{Say it:} \emph{suvidhā}, once more
  \item \emph{Say it:} \emph{yā viśrāmgṛhāt sagaḷyā suvidhā āhet.} — This guest house has all the amenities.
  \item \emph{Recall:} say the Marathi for to ask after, then the Marathi for to drop by, then say \emph{suvidhā} again
\end{itemize}

\begin{tcolorbox}[breakable,colback=teal!4,colframe=teal!35!black,title={Before you move on}]
What does \mr{सुविधा} mean? (\textquotedblleft{}An amenity\textquotedblright{}.) Say it once more.
\end{tcolorbox}
\section[\texorpdfstring{\mr{बिछाना} (bichānā)}{बिछाना (bichānā)}]{\mr{बिछाना} (bichānā) — bedding; a bed}

\label{lesson:MR-C415-bichana}



\begin{quote}
Before the new one: say the Marathi for to drop by, then the Marathi for an amenity.
\end{quote}

\subsection*{You'll want to know: \mr{बिछाना}}
\textbf{\mr{बिछाना}} — \emph{bichānā} — \textquotedblleft{}bedding; a bed\textquotedblright{}.

\begin{quote}
\mr{आईने} \mr{पाहुण्यांसाठी} \mr{बिछाना} \mr{तयार} \mr{केला}.
\end{quote}

\begin{quote}
\emph{āīne pāhuṇyāṁsāṭhī bichānā tayār kelā.}
\end{quote}

\begin{quote}
Mother got the bedding ready for the guests.
\end{quote}

Masculine.

The mattress, sheet and pillow together, laid out for sleep.

\begin{cousinweb}
\textbf{\mr{बिछाना}} (\emph{bichānā}) is a loan, usually traced to Hindustani: CDIAL lists Hindi \emph{bichānā}, \textquotedblleft{}to spread\textquotedblright{}, under \emph{*vicchādayati}, \textquotedblleft{}spread out\textquotedblright{} (CDIAL 11692).
\end{cousinweb}

\begin{grammarlens}[title={What you've built}]
One more word for getting around.
\end{grammarlens}

\subsection*{Your turn}
\begin{itemize}
\raggedright
  \item \emph{Say it:} \emph{bichānā}
  \item \emph{Say it:} \emph{bichānā}, once more
  \item \emph{Say it:} \emph{āīne pāhuṇyāṁsāṭhī bichānā tayār kelā.} — Mother got the bedding ready for the guests.
  \item \emph{Recall:} say the Marathi for to drop by, then the Marathi for an amenity, then say \emph{bichānā} again
\end{itemize}

\begin{tcolorbox}[breakable,colback=teal!4,colframe=teal!35!black,title={Before you move on}]
What does \mr{बिछाना} mean? (\textquotedblleft{}Bedding; a bed\textquotedblright{}.) Say it once more.
\end{tcolorbox}
\section[\texorpdfstring{\mr{घोंगडी} (ghoṅgḍī)}{घोंगडी (ghoṅgḍī)}]{\mr{घोंगडी} (ghoṅgḍī) — a woollen blanket}

\label{lesson:MR-C415-ghongdi}



\begin{quote}
Before the new one: say the Marathi for an amenity, then the Marathi for bedding.
\end{quote}

\subsection*{You'll want to know: \mr{घोंगडी}}
\textbf{\mr{घोंगडी}} — \emph{ghoṅgḍī} — \textquotedblleft{}a woollen blanket\textquotedblright{}.

\begin{quote}
\mr{थंडीत} \mr{आजोबा} \mr{घोंगडी} \mr{अंगावर} \mr{घेतात}.
\end{quote}

\begin{quote}
\emph{thaṇḍīt ājobā ghoṅgḍī aṅgāvar ghetāt.}
\end{quote}

\begin{quote}
In the cold Grandfather wraps himself in a woollen blanket.
\end{quote}

Feminine.

A rough blanket of sheep’s wool, woven by shepherds.

\begin{cousinweb}
\textbf{\mr{घोंगडी}} (\emph{ghoṅgḍī}) is a Marathi word for the coarse wool blanket shepherds weave; its origin is not settled.
\end{cousinweb}

\begin{grammarlens}[title={What you've built}]
One more word for getting around.
\end{grammarlens}

\subsection*{Your turn}
\begin{itemize}
\raggedright
  \item \emph{Say it:} \emph{ghoṅgḍī}
  \item \emph{Say it:} \emph{ghoṅgḍī}, once more
  \item \emph{Say it:} \emph{thaṇḍīt ājobā ghoṅgḍī aṅgāvar ghetāt.} — In the cold Grandfather wraps himself in a woollen blanket.
  \item \emph{Recall:} say the Marathi for an amenity, then the Marathi for bedding, then say \emph{ghoṅgḍī} again
\end{itemize}

\begin{tcolorbox}[breakable,colback=teal!4,colframe=teal!35!black,title={Before you move on}]
What does \mr{घोंगडी} mean? (\textquotedblleft{}A woollen blanket\textquotedblright{}.) Say it once more.
\end{tcolorbox}
\section[\texorpdfstring{\mr{अंथरणे} (antharṇe)}{अंथरणे (antharṇe)}]{\mr{अंथरणे} (antharṇe) — to spread out bedding}

\label{lesson:MR-C415-antharne}



\begin{quote}
Before the new one: say the Marathi for bedding, then the Marathi for a woollen blanket.
\end{quote}

\subsection*{You'll want to know: \mr{अंथरणे}}
\textbf{\mr{अंथरणे}} — \emph{antharṇe} — \textquotedblleft{}to spread out bedding\textquotedblright{}.

\begin{quote}
\mr{आईने} \mr{जमिनीवर} \mr{चादर} \mr{अंथरली}.
\end{quote}

\begin{quote}
\emph{āīne jaminīvar cādar antharlī.}
\end{quote}

\begin{quote}
Mother spread a sheet out on the floor.
\end{quote}

Past, with -\mr{ने} (he …): \mr{त्याने} \mr{अंथरलं}, \emph{tyāne antharlaṁ}.

The thing spread is the object: \mr{चादर} \mr{अंथरली}, \mr{गादी} \mr{अंथरली}.

\begin{cousinweb}
\textbf{\mr{अंथरणे}} (\emph{antharṇe}) is inherited: CDIAL lists Marathi \emph{ā̃tharṇẽ}, \textquotedblleft{}to spread out\textquotedblright{}, under Sanskrit \emph{āstarati}, \textquotedblleft{}spreads\textquotedblright{} (CDIAL 1507).
\end{cousinweb}

\begin{grammarlens}[title={What you've built}]
One more word for getting around.
\end{grammarlens}

\subsection*{Your turn}
\begin{itemize}
\raggedright
  \item \emph{Say it:} \emph{antharṇe}
  \item \emph{Say it:} \emph{antharṇe}, once more
  \item \emph{Say it:} \emph{āīne jaminīvar cādar antharlī.} — Mother spread a sheet out on the floor.
  \item \emph{Recall:} say the Marathi for bedding, then the Marathi for a woollen blanket, then say \emph{antharṇe} again
\end{itemize}

\begin{tcolorbox}[breakable,colback=teal!4,colframe=teal!35!black,title={Before you move on}]
What does \mr{अंथरणे} mean? (\textquotedblleft{}To spread out bedding\textquotedblright{}.) Say it once more.
\end{tcolorbox}
\section[\texorpdfstring{\mr{पांघरणे} (pāṅgharṇe)}{पांघरणे (pāṅgharṇe)}]{\mr{पांघरणे} (pāṅgharṇe) — to cover oneself}

\label{lesson:MR-C415-pangharne}



\begin{quote}
Before the new one: say the Marathi for a woollen blanket, then the Marathi for to spread out bedding.
\end{quote}

\subsection*{You'll want to know: \mr{पांघरणे}}
\textbf{\mr{पांघरणे}} — \emph{pāṅgharṇe} — \textquotedblleft{}to cover oneself\textquotedblright{}.

\begin{quote}
\mr{रात्री} \mr{मी} \mr{जाड} \mr{चादर} \mr{पांघरली}.
\end{quote}

\begin{quote}
\emph{rātrī mī jāḍ cādar pāṅgharlī.}
\end{quote}

\begin{quote}
At night I covered myself with a thick sheet.
\end{quote}

Past, with -\mr{ने} (he …): \mr{त्याने} \mr{पांघरलं}, \emph{tyāne pāṅgharlaṁ}.

The cover is the object, and the verb agrees with it: \mr{चादर} \mr{पांघरली}.

\begin{cousinweb}
\textbf{\mr{पांघरणे}} (\emph{pāṅgharṇe}) is a Marathi verb for pulling a cover over yourself; its origin is not settled; the noun \textbf{\mr{पांघरूण}}, a covering, belongs with it.
\end{cousinweb}

\begin{grammarlens}[title={What you've built}]
One more word for getting around.
\end{grammarlens}

\subsection*{Your turn}
\begin{itemize}
\raggedright
  \item \emph{Say it:} \emph{pāṅgharṇe}
  \item \emph{Say it:} \emph{pāṅgharṇe}, once more
  \item \emph{Say it:} \emph{rātrī mī jāḍ cādar pāṅgharlī.} — At night I covered myself with a thick sheet.
  \item \emph{Recall:} say the Marathi for a woollen blanket, then the Marathi for to spread out bedding, then say \emph{pāṅgharṇe} again
\end{itemize}

\begin{tcolorbox}[breakable,colback=teal!4,colframe=teal!35!black,title={Before you move on}]
What does \mr{पांघरणे} mean? (\textquotedblleft{}To cover oneself\textquotedblright{}.) Say it once more.
\end{tcolorbox}
